\documentclass[11pt]{article}
\usepackage[T1]{fontenc}
\usepackage{lmodern,microtype}
\usepackage{amsmath,amssymb,amsthm,mathtools}
\usepackage[letterpaper,margin=1in]{geometry}
\usepackage{booktabs,array,enumitem,longtable}
\usepackage{xcolor}
\usepackage[colorlinks=true,linkcolor=blue!55!black,citecolor=blue!55!black,urlcolor=blue!55!black]{hyperref}
\newcommand{\F}{F}
\newcommand{\K}{K}
\newcommand{\A}{\mathbb A}
\newcommand{\C}{\mathbb C}
\newcommand{\Z}{\mathbb Z}
\newcommand{\GL}{\operatorname{GL}}
\newcommand{\Irr}{\operatorname{Irr}}
\newcommand{\Nrd}{\operatorname{Nrd}}
\newcommand{\Nm}{\operatorname{N}}
\newcommand{\vol}{\operatorname{vol}}
\newcommand{\Gal}{\operatorname{Gal}}
\newcommand{\Aut}{\operatorname{Aut}}
\newcommand{\reg}{\mathrm{reg}}
\newcommand{\tH}{\widetilde H}
\newcommand{\tZ}{\widetilde Z}
\newcommand{\tT}{\widetilde T}
\newcommand{\Gb}{\overline G}
\newcommand{\GDb}{\overline{G_D}}
\newcommand{\Ptag}{\textnormal{\textbf{[P]}}}
\newcommand{\Etag}{\textnormal{\textbf{[E]}}}
\newcommand{\Ctag}{\textnormal{\textbf{[C]}}}
\newcommand{\Otag}{\textnormal{\textbf{[O]}}}
\theoremstyle{plain}
\newtheorem{theorem}{Theorem}[section]
\newtheorem{proposition}[theorem]{Proposition}
\newtheorem{lemma}[theorem]{Lemma}
\newtheorem{corollary}[theorem]{Corollary}
\theoremstyle{definition}
\newtheorem{definition}[theorem]{Definition}
\newtheorem{remark}[theorem]{Remark}
\newtheorem{example}[theorem]{Example}
\title{A metaplectic Jacquet--Langlands transfer for division algebras of prime degree, IV:\\
Wild local correspondence and global trace formula comparison\\[0.4em]\normalsize Draft 2}
\author{David Joyner \and ChatGPT (OpenAI) \and Claude (Anthropic)}
\date{November 7, 2026}
\emergencystretch=3em
\begin{document}
\maketitle
\begin{abstract}
We continue Parts I--III for an odd-degree covering of the multiplicative
inner form of $\GL_r$ with $r$ prime and $(r,n)=1$.  First we reconcile the
split-diagonal cocycle with the ambient-cover restriction and central
relabeling in Part I.  The regular good-element criterion for \emph{every}
maximal torus, including nonsplit tori and wild residue characteristic,
is already a theorem of Flicker--Kazhdan; we supply its proof with the
parameters used in this series.  We prove two refinements useful in global
comparison: an almost-everywhere local-to-global criterion for globally
good regular classes, and an exact local calculation of the kernel, index,
and Haar modulus of the $n$th-power map on an arbitrary finite \emph{\'etale}
torus modulo scalars.  The latter makes the split-place root multiplicity
explicit and shows why regularity of a root must be distinguished from
regularity of its $n$th power.  The genuine character transform of Part III
is then placed alongside the corrected metaplectic character relation of
Zou and the older work of Flicker--Kazhdan.  A normalization ledger and
geometric--spectral proof obligations separate the previously established
covered-inner-form trace comparison from the additional proposed
comparison with ordinary $D^\times$.  The wild irreducible-character
integrality and the latter global comparison are not established here.
\end{abstract}
\tableofcontents

\section{Introduction and mathematical status}\label{sec:intro}
The problem concerns three different local representation theories.
Write $H_v=D_v^\times$, and let $\GDb_v$ denote its specified $n$-fold
central cover; the cover $\Gb_v$ of $\GL_r(K_v)$ is defined inside the
same ambient cover.  An additional \emph{ordinary} group $H_v$ enters
through the transform $L_{D,v}$ of Part III.  These are not interchangeable.
In particular, a comparison between $\GDb$ and $\Gb$ has already been
obtained in Parts I--II under its stated test-function and spectral
hypotheses.  The new proposed comparison between $H$ and $\GDb$ is a
different assertion.

\paragraph{Global assumptions.}
Let $K$ be a totally imaginary number field with $\mu_n\subset K$,
$n\geq1$ odd.  Let $r\geq2$ be prime, $(r,n)=1$, let $L/K$ be cyclic
of degree $r$, and let $D$ be a central division $K$-algebra of
degree $r$ containing $L$.  Put $S=\{v:D_v\text{ is division}\}$.
Fix integers $r',c$ such that
\begin{equation}\label{eq:parameters}
 rr'\equiv1\pmod n,\qquad r-1+2c\equiv0\pmod n,
 \qquad c_v\equiv r'c\pmod n.
\end{equation}
Thus $r-1+2rc_v\equiv0\pmod n$.  In fact $c_v$ is independent
of $v$ with these choices; the subscript recalls Part I's notation.
Locally $F=K_v$ is a nonarchimedean field containing $\mu_n$;
$p$ denotes its residue characteristic.  We call $v$ \emph{wild}
if $p\mid n$.

\paragraph{Status convention.}
\textbf{[P]} denotes a proof given in this paper (including cited
inputs), \textbf{[E]} a result established in a cited publication or an
earlier part, \textbf{[C]} an implication conditional on specified
unproved hypotheses, and \textbf{[O]} an open objective.  No
representation-theoretic correspondence follows from an $L^2$ isometry
without further input.

\begin{center}\small
\begin{tabular}{@{}p{0.49\linewidth}p{0.37\linewidth}p{0.07\linewidth}@{}}
\toprule
Assertion & Precise basis & Status\\\midrule
Diagonal and nonsplit good-element criterion, including wild $p$ & [FK, \S3]; Part I, \S3 & [E]\\
Compatibility of the Part I cocycle with $\sigma_{c_v}$ & Part I, Prop. 2.3(b); direct calculation & [P]\\
Power-map kernel, index, and Haar modulus on $T_E(F)/F^\times$ & Section \ref{sec:localpowers} & [P]\\
Global good-regular class criterion & Section \ref{sec:globalgood}; Part III Hasse principle & [P]\\
Division-side unitary character transform and tame integral case & Part III, Thms. 4.3, 7.6 & [E]\\
Wild irreducible-character integrality & Section \ref{sec:wild} & [O]\\
Covered $D$--covered $\GL$ global simple trace comparison & Parts I--II under their hypotheses & [E]\\
Covered $D$--ordinary $D$ global comparison & Section \ref{sec:geometric} & [O]\\
Removal of auxiliary support and packet hypotheses & Section \ref{sec:spectral} & [O]\\
\bottomrule
\end{tabular}
\end{center}

\paragraph{Scope of direct cross-checking.}
The cocycle and center calculations below have been checked against
Part I, Draft 5, especially its Propositions 2.3 and 3.1, and the
original [FK, \S\S2--4].  Claude's attached referee report independently
compared these formulas with Part I, Draft 6, and Parts II--III.
We also used Part II, Draft 4, and Part III, Draft 7 to check the
relative normalizations.  This provenance identifies exactly which
versions were available for direct verification; it should not be read
as a claim of a line-by-line independent re-refereeing of every theorem
in Parts I--III.

\section{The exact cocycle convention and star map}\label{sec:cocycle}
This section fixes the one place where a parameter mismatch could
invalidate every subsequent constant.  The ambient construction is
a restriction of a cocycle over $L_v$, whereas [FK] uses the standard
$\GL_r(F)$ Kazhdan--Patterson convention.

\subsection{Restriction and relabeling}
Let $\sigma_F$ be the parameter-zero Kazhdan--Patterson cocycle with
restriction to the diagonal torus
\begin{equation}\label{eq:parameterzero}
\sigma_F(\operatorname{diag}(a_i),\operatorname{diag}(b_i))
=\prod_{i<j}(a_i,b_j)_{n,F}.
\end{equation}
The $\GL$ subgroup of the ambient cover in Part I has, \emph{as an
extension}, cocycle model
\begin{equation}\label{eq:before}
  \sigma_F(g,h)^r\,(\det g,\det h)_{n,F}^{c}
  \quad\text{before relabeling }\mu_n.
\end{equation}
This assertion is Part I, Proposition 2.3(b).  Applying the
kernel relabeling $\zeta\mapsto\zeta^{r'}$ yields the model
\begin{equation}\label{eq:after}
 \sigma_{c_v}(g,h)=\sigma_F(g,h)
 (\det g,\det h)_{n,F}^{c_v}.
\end{equation}
On the diagonal torus, this is exactly the cocycle used in
Part IV, Draft 1; its letter $c$ meant the \emph{new} $c_v$,
not the $c$ of \eqref{eq:parameters}.  The distinction is essential.
The isomorphism in \eqref{eq:before} is an isomorphism of extensions,
not a claim that arbitrary choices of cocycle sections give
literally equal functions on all of $\GL_r(F)^2$.

\begin{proposition}[\Ptag\ Diagonal commutator audit]\label{prop:diag}
For $t=\operatorname{diag}(t_1,\ldots,t_r)$ and
$u=\operatorname{diag}(u_1,\ldots,u_r)$ in $\GL_r(F)$,
\begin{equation}\label{eq:matrix}
 [\widetilde t,\widetilde u]
 =\prod_{i,j}(t_i,u_j)_{n,F}^{A_{ij}},\qquad
 A=-I_r+(1+2c_v)J_r.
\end{equation}
Here $J_r$ is the all-ones matrix.  Its row sums are
$r-1+2rc_v\equiv0\pmod n$, so scalar lifts commute with the lifted
split torus.  The radical on $(F^\times/F^{\times n})^r$ consists
precisely of diagonal classes.
\end{proposition}
\begin{proof}
The commutator is the ratio
$\sigma_{c_v}(t,u)/\sigma_{c_v}(u,t)$.
The Hilbert symbol satisfies $(a,b)_n=(b,a)_n^{-1}$; this is
\emph{antisymmetry}, valid for all $n$.  The determinant twist
contributes $2c_v$ to each entry of $A$, and the parameter-zero
factor contributes $1$ off the diagonal.  This gives \eqref{eq:matrix}.
The Hilbert pairing on $X=F^\times/F^{\times n}$ is perfect.
As $X$ is killed by $n$ and $(r,n)=1$, the decomposition
$X^r=\Delta X\oplus\{(x_i):\sum x_i=0\}$ holds by averaging with
$r^{-1}$ modulo $n$.  On these summands $A$ acts, respectively, by
$r-1+2rc_v=0$ and $-1$.  Therefore $\ker A=\Delta X$.
\end{proof}

\begin{remark}
The assumption $(r,n)=1$ also follows algebraically from
$n\mid r-1+2rc_v$, because any common divisor of $r$ and $n$ must
divide $1$.  We retain it because it is a global hypothesis of the
series.  The oddness of $n$ is used for the ordinary star-map
convention below, not for Hilbert-symbol antisymmetry.
\end{remark}

\subsection{Star and central characters}\label{subsec:star}
For an element $h$ of the ordinary group, choose any lift
$\widetilde h$ and define $h^*=\widetilde h^{\,n}$.  It is
independent of the lift because $\mu_n$ is central.  In the present
scalar-central cover, for $a\in F^\times$,
\begin{equation}\label{eq:starcentral}
 (ah)^*=a^*h^*,\qquad a^*=\widetilde a^{\,n}.
\end{equation}
If $h_1,h_2$ commute in an ordinary torus, their lifted commutator
is in $\mu_n$, and the class-two power identity
\begin{equation}\label{eq:starpower}
 (\widetilde h_1\widetilde h_2)^n
 =\widetilde h_1^{\,n}\widetilde h_2^{\,n}
 [\widetilde h_2,\widetilde h_1]^{n(n-1)/2}
 =\widetilde h_1^{\,n}\widetilde h_2^{\,n}
\end{equation}
uses $n$ odd.  Thus star is multiplicative on the lifted torus.
This agrees with the convention in [FK, \S4] and [Zou, \S3]:
Zou writes $x^*=u(x)s(x)^n$, but for odd $n$ the correction
$u(x)$ is identically $1$.  In particular there is no hidden
sign in our odd-degree star map.  For a scalar $\zeta\in\mu_n(F)$,
$\zeta^*=1$ in the standard diagonal cocycle model: the
self-symbol $(\zeta,\zeta)_n=1$ for odd $n$, so
$s(\zeta I_r)^n=s(I_r)=1$.  This last statement is also
invariant under the extension isomorphisms in Part I.
For representation characters we first fix a genuine unitary
character $\widetilde\theta$ of the central preimage $\tZ$ and
\emph{then} put $\theta(a)=\widetilde\theta(a^*)$.
The reverse implication is not unique: $\theta$ alone does not
specify $\widetilde\theta$.  Test functions use the inverse
central characters $\theta^{-1}$ and $\widetilde\theta^{-1}$.

\section{Good elements for every maximal torus}\label{sec:good}
A regular maximal torus in $\GL_r(F)$ may be nonsplit even at a
place where $D_v\simeq M_r(F)$.  This is the situation not treated by
the diagonal matrix alone.  The following calculation is not new:
it is the proof of [FK, \S3, Proposition] in the notation of this
series, and also underlies Part I, Proposition 3.5.

\begin{theorem}[\Etag\ Flicker--Kazhdan good-element theorem]\label{thm:FKgood}
Let $E=\prod_{i=1}^kE_i$ be a finite \emph{\'etale} $F$-algebra
of dimension $r$, embedded as a regular maximal torus
$T_E(F)=E^\times$ in $\GL_r(F)$.  Under \eqref{eq:parameters},
for any $t\in T_E(F)$ and any lift $\widetilde t$,
\begin{equation}\label{eq:FKtorus}
 [\widetilde t,\widetilde u]
 =(\Nm_{E/F}t,\Nm_{E/F}u)_{n,F}^{1+2c_v}
   \prod_{i=1}^k(t_i,u_i)_{n,E_i}^{-1}.
\end{equation}
The radical of this commutator pairing is
\begin{equation}\label{eq:everygood}
 \{t\in T_E(F):[\widetilde t,\widetilde u]=1\ \forall u\in T_E(F)\}
 =F^\times T_E(F)^n.
\end{equation}
For regular $t$, this is exactly the criterion that $\widetilde t$
is good.  Both identities are valid when $p\mid n$.
\end{theorem}
\begin{proof}[Proof from the cited formula, included for clarity]
The commutator formula is [FK, \S3], using [KP, Prop.~0.1.5];
its restriction to a single field is exactly Part I,
Proposition 3.1 after relabeling.  Write $N=\Nm_{E/F}$.
Norm compatibility of local Hilbert symbols gives
\[
 (Nt,Nu)_{n,F}
 =\prod_i(Nt,u_i)_{n,E_i}.
\]
Thus \eqref{eq:FKtorus} equals
\[
 \prod_i\bigl((Nt)^{1+2c_v}t_i^{-1},u_i\bigr)_{n,E_i}.
\]
The nondegeneracy of the $n$th Hilbert symbol over each $E_i$
(including wildly ramified $E_i$) implies that the commutator is
trivial for all $u$ exactly when
\begin{equation}\label{eq:pointwisegood}
 t_i\equiv(Nt)^{1+2c_v}\pmod{E_i^{\times n}}
 \quad (1\leq i\leq k).
\end{equation}
This directly says $t\in F^\times E^{\times n}$.  Conversely,
if $t=ay^n$ for $a\in F^\times$, then
$Nt=a^r(Ny)^n$; therefore
$(Nt)^{1+2c_v}\equiv a^{r(1+2c_v)}\equiv a\equiv t_i$
modulo $E_i^{\times n}$ by scalar centrality.  This proves the
criterion.  No tameness assumption has been introduced.
\end{proof}

\begin{example}[The diagonal rank-two check]\label{ex:diag}
Take $r=2$, $n=3$ and $F=\mathbb Q_3(\zeta_3)$.
In the notation of Part I, Example 3.7, choose $c=1$, so
$c_v\equiv2\pmod3$.  Formula \eqref{eq:matrix} becomes
\[
 A\equiv\begin{pmatrix}1&-1\\-1&1\end{pmatrix}\pmod3,
 \qquad
 [\widetilde t,\widetilde u]
 =(t_1/t_2,u_1/u_2)_{3,F}.
\]
Thus $\operatorname{diag}(t_1,t_2)$, when regular, is good exactly
when $t_1/t_2\in F^{\times3}$.  If a quadratic field $E/F$ is
embedded as a nonsplit maximal torus, the same theorem says that
goodness is equivalent to $t\in F^\times E^{\times3}$; the two
tests are instances of the same formula, not separate conjectures.
\end{example}

\begin{remark}[Why goodness is not enough for a regular-root formula]\label{rem:power-singular}
The power of a regular element in a \emph{split} torus need not be
regular.  For $n=3$ the element
$\operatorname{diag}(1,\zeta_3)$ is regular, but its cube is the
identity.  A character or orbital-integral formula with
$\Delta(t^n)$ in the denominator must restrict to the set where
$t^n$ is regular (as Zou explicitly does), or use a separately
proved germ extension.  On a degree-$r$ \emph{field} torus with
$(r,n)=1$ and $r$ prime, regularity is preserved by $n$th powers.
\end{remark}

\section{A new local power-map calculation for \texorpdfstring{\'etale}{etale} tori}\label{sec:localpowers}
The key local distinction is between a field centralizer (compact
modulo scalars) and a product of fields (not compact modulo scalars).
Both occur in the split \emph{group}.  The next proposition gives
an exact count of the roots and the size of the image at wild and
tame places, including the factors that a global trace identity
must account for.

\begin{lemma}[\Ptag\ Kummer injection for a product algebra]\label{lem:Kumprod}
Let $E=\prod_{i=1}^kE_i$ be finite \emph{\'etale} of dimension $r$
over $F$, with $\mu_n\subset F$ and $(r,n)=1$.  Then
\begin{equation}\label{eq:Kummerinj}
 F^\times/F^{\times n}\longrightarrow
 E^\times/E^{\times n}
\end{equation}
is injective.
\end{lemma}
\begin{proof}
Suppose $a\in F^\times$ is an $n$th power in every $E_i$.
The Kummer extension $F(a^{1/n})/F$, since $\mu_n\subset F$,
is cyclic of degree $d\mid n$ and embeds in every $E_i$.
Therefore $d\mid[E_i:F]$ for every $i$, so also
$d\mid\sum_i[E_i:F]=r$.  Coprimality gives $d=1$ and
$a\in F^{\times n}$.
\end{proof}

\begin{proposition}[\Ptag\ Kernel, index, and Haar modulus]\label{prop:power-etale}
Under the hypotheses of Lemma \ref{lem:Kumprod}, put
$A_E=E^\times/F^\times$ and $P_n([t])=[t^n]$.  Then
\begin{align}
 \ker(P_n)&\simeq \mu_n(E)/\mu_n(F)
      \simeq \mu_n^{\,k}/\mu_n,
 &|\ker(P_n)|&=n^{k-1},\label{eq:kernel}\\
 [A_E:P_n(A_E)]&=n^{2(k-1)}|n|_F^{-(r-1)}.
 \label{eq:index}
\end{align}
In particular, $P_n(A_E)$ is open and closed.  With \emph{any}
Haar measure $da$ on $A_E$ and its restriction to $P_n(A_E)$,
for compactly supported integrable $q$ on that image,
\begin{equation}\label{eq:haarJac}
 \int_{A_E}q(a^n)\,da
 =n^{k-1}|n|_F^{-(r-1)}
   \int_{P_n(A_E)}q(b)\,db.
\end{equation}
\end{proposition}
\begin{proof}
If $t^n=a\in F^\times$, then $a$ is an $n$th power in every
$E_i$; by Lemma \ref{lem:Kumprod}, write $a=b^n$ with $b\in F^\times$.
Thus each $t_i/b\in\mu_n(E_i)=\mu_n(F)$, giving
\eqref{eq:kernel}, including its canonical description.

For a local field $M$ containing $\mu_n$ the Kummer quotient has
\begin{equation}\label{eq:localKumIndex}
 |M^\times/M^{\times n}|=n^2|n|_M^{-1}.
\end{equation}
Indeed the discrete valuation contributes $n$; the finite
torsion unit factor containing $\mu_n$ contributes $n$; and
the torsion-free $p$-adic unit factor contributes
$|n|_M^{-1}$ (multiplication by $n$ on its $\Z_p$-lattice).  Alternatively, \eqref{eq:localKumIndex}
follows from local duality.
Since $F^\times/F^{\times n}$ injects into
$E^\times/E^{\times n}$, the group $A_E/A_E^n$ is its cokernel:
\[
 |A_E/A_E^n|
 =\frac{\prod_i n^2|n|_{E_i}^{-1}}
             {n^2|n|_F^{-1}}
 =n^{2(k-1)}|n|_F^{-(r-1)},
\]
where $|n|_{E_i}=|n|_F^{[E_i:F]}$ for normalized absolute values
on local fields.  This proves \eqref{eq:index}.

For the Haar statement, the open endomorphism of this locally compact
abelian group with finite kernel $K$ and open finite-index image
$I$ pushes Haar measure to a Haar measure on its image.  Its density
relative to Haar measure restricted from the ambient group equals
$[A_E:I]/|K|$ for this power map; this can be checked explicitly
on $\mathbb Z^{k-1}$ (counting measure) and a compact open subgroup
(probability Haar), using that multiplication by $n$ has valuation
index $n^{k-1}$.  Substitution of \eqref{eq:kernel} and
\eqref{eq:index} gives \eqref{eq:haarJac}.
\end{proof}

\begin{remark}[A measure-theoretic warning]\label{rem:haar}
One should not confuse the \emph{index} of the image with the
Jacobian factor in \eqref{eq:haarJac}.  When the group is noncompact,
the kernel and the free valuation lattice contribute differently.
The scalar in \eqref{eq:haarJac} is the ratio
$[A_E:P_nA_E]/|\ker P_n|$, not simply the index.
\end{remark}

\begin{example}[Three local torus types]
If $E/F$ is a field of degree $r$, then $k=1$ and $P_n$ is
injective, with index and Haar modulus $|n|_F^{-(r-1)}$.
If $E=F^r$ then $k=r$ and the kernel has order $n^{r-1}$.
For $r=2,n=3$, $F=\mathbb Q_3(\zeta_3)$ and $|3|_F=1/9$,
\[
\begin{array}{c|ccc}
\text{type}&|\ker P_3|&[A_E:A_E^3]&\text{factor in \eqref{eq:haarJac}}\\\hline
E/F\text{ quadratic}&1&9&9\\
E=F\times F&3&81&27
\end{array}
\]
These integers quantify why one cannot use the division-torus
substitution formula unchanged on the full split group.
\end{example}

\section{A global classification of good regular classes}\label{sec:globalgood}
The local theorem now applies to the actual torus at \emph{every}
place.  Combining it with the Hasse principle of Part III removes
the local-to-global ambiguity for power classes of a degree-$r$
field, though not the independent measure and spectral questions.

\begin{proposition}[\Etag\ Almost-everywhere powers modulo scalars]\label{prop:Hasse}
Let $E/K$ be a field extension of degree $r$ with $(r,n)=1$ and
$\mu_n\subset K$.  For every finite set $\Sigma$ of places,
\begin{equation}\label{eq:Hasse}
 E^\times/(K^\times E^{\times n})\longrightarrow
 \prod_{v\notin\Sigma}
 \frac{(E\otimes_K K_v)^\times}
 {K_v^\times(E\otimes_KK_v)^{\times n}}
\end{equation}
is injective.  This is Part III, Proposition 11.3.
\end{proposition}
\begin{proof}[Recalled proof]
Let $t$ be locally scalar times an $n$th power outside $\Sigma$.
Choose $u,k\in\Z$ with $ru+nk=1$ and put
$a=N_{E/K}(t)^u\in K^\times$.  If $t=a_v y_v^n$ in the local
\'etale algebra, taking norms gives the identity
\[
 t/a=\bigl(a_v^k y_v/N_{E_v/K_v}(y_v)^u\bigr)^n.
\]
Thus $t/a$ is an $n$th power in every completion of $E$ over all
$v\notin\Sigma$.  If its class in $E^\times/E^{\times n}$ were
nontrivial, the nontrivial Kummer extension
$E((t/a)^{1/n})/E$ would, by Chebotarev, admit infinitely many
unramified places at which $t/a$ is not an $n$th power.
This contradicts the hypothesis.  Hence $t\in K^\times E^{\times n}$.
\end{proof}

\begin{theorem}[\Ptag\ Global good-regular criterion]\label{thm:globalgood}
Assume the covers and rational splittings of Part I.  Let
$G$ denote $D^\times$ or $\GL_r$, and let $\gamma\in G(K)$ be
regular elliptic, so $E=K[\gamma]$ is a field of degree $r$.
The following are equivalent:
\begin{enumerate}[label=(\alph*)]
\item the rational lift $\iota(\gamma)$ is good at every place;
\item it is good at every place outside some finite set $\Sigma$;
\item $\gamma\in K^\times E^{\times n}$.
\end{enumerate}
The map $[s]\mapsto[s^n]$ is a bijection from the noncentral
classes in $E^\times/K^\times$ onto the globally good regular
classes belonging to $E$.  Furthermore
\begin{equation}\label{eq:Wstabilizer}
 W_{s^n}=W_s,\qquad
 W_s=\{\sigma\in\Aut(E/K):\sigma(s)\in K^\times s\}.
\end{equation}
\end{theorem}
\begin{proof}
At each finite $v$, the ordinary centralizer is
$T_v=(E\otimes_K K_v)^\times$.  If $D_v$ splits, apply
Theorem \ref{thm:FKgood}, transported through the exact Part I
cover isomorphism.  If $D_v$ is division, use Part I,
Proposition 3.5(a), whose commutator is the same by Part I,
Proposition 3.1.  At complex places the $n$th-power map on
$T_v$ is surjective.  Thus at every place the local goodness
condition is precisely
$\gamma\in K_v^\times T_v^n$.
Implications (c)$\Rightarrow$(a)$\Rightarrow$(b) follow,
while (b)$\Rightarrow$(c) is Proposition \ref{prop:Hasse}.

The power map $E^\times/K^\times\to E^\times/K^\times$
is injective: if $s^n\in K^\times$, Kummer theory shows
$[K(s):K]\mid n$, while $[K(s):K]\mid [E:K]=r$.
Hence $s\in K^\times$.  Because $r$ is prime, every noncentral
$s\in E$ generates $E$, and the same is true of $s^n$.
This proves the bijection.  Finally
$\sigma(s^n)/s^n\in K^\times$ if and only if
$(\sigma(s)/s)^n\in K^\times$, which is equivalent by injectivity
to $\sigma(s)/s\in K^\times$.  Equation \eqref{eq:Wstabilizer}
follows.
\end{proof}

\begin{remark}[Precise consequence for the geometric side]
Part I, Proposition 6.1, uses the rational regular classes which
are good everywhere, with coefficient
\begin{equation}\label{eq:globalvol}
 c(\gamma)=\frac{\vol\bigl(\A_K^\times E^\times\backslash
 \A_E^\times;dt\bigr)}{|W_\gamma|}.
\end{equation}
Theorem \ref{thm:globalgood} parametrizes its noncentral
summation set by global $[s^n]$ and preserves $W_\gamma$.
It does \emph{not} replace the local root sums in Zou's
character formula: the global root class may be unique even when
a localized split torus has $n^{k_v-1}$ roots modulo scalars.
\end{remark}

\section{The remaining wild local character problem}\label{sec:wild}
We distinguish character transfer from class-function transfer.
The latter is a theorem of Part III; the former requires an
\emph{integral} statement not implied by unitarity.

\subsection{Division-side normalization and discriminants}
Let $F$ be nonarchimedean, $D/F$ a division algebra of prime
degree $r$, $H=D^\times$, and $\tH$ its specified cover.
Set $Z=F^\times$ and $\tZ=p^{-1}(Z)$, equipped with compatible
Haar measures on $H/Z\simeq\tH/\tZ$.  The Part III transform is
\begin{equation}\label{eq:LD}
 (L_D\Phi)(z\,t^*)
 =b_D\,\widetilde\theta(z)
  \frac{\Delta(t)}{\Delta(t^n)}\Phi(t),\qquad
 b_D=|n|_F^{-(r-1)/2},
\end{equation}
for regular $t$, $z\in\tZ$; it vanishes at nongood regular lifts.
Here
\begin{equation}\label{eq:Delta}
 \Delta(t)=\left|\det\bigl(1-\operatorname{Ad}(t)
 \mid\mathfrak h/\mathfrak t\bigr)\right|_F^{1/2}
 =\prod_{i<j}\left|
 \frac{(\lambda_i-\lambda_j)^2}{\lambda_i\lambda_j}
 \right|_F^{1/2}.
\end{equation}
This equals Part I's $|D(t)|^{1/2}|\Nrd(t)|_F^{-(r-1)/2}$.
Since the power map on a field torus modulo scalars has index
$|n|_F^{-(r-1)}$, Weyl integration gives $b_D^2$
equal to this index.  The factor $b_D$ is independent of
\emph{common} rescalings of the two identified quotient Haar
measures.  The ratio of discriminants extends at central elements
to $|n|_F^{-r(r-1)/2}$ (Part III, Lemma 3.3), hence
\begin{equation}\label{eq:cD}
 (L_D\Theta_\rho)(1)=c_D\dim\rho,\qquad
 c_D=b_D|n|_F^{-r(r-1)/2}
 =|n|_F^{-(r^2-1)/2}.
\end{equation}
The identity factor is pointwise and does not depend on Haar
measures.  Part III, Theorem 9.1, separately gives the adjoint
orbital-integral/test-function transfer.  Test functions transform
by $\widetilde\theta^{-1}$, not by $\widetilde\theta$.

\subsection{The integral matrix}\label{subsec:integral}
Choose $\widetilde\theta$ first and put
$\theta(a)=\widetilde\theta(a^*)$.
Part III, Theorem 4.3 and Corollary 5.1, show that $L_D$ is a
unitary \emph{isomorphism} of invariant $L^2$ spaces and that
\begin{equation}\label{eq:matrixchars}
 L_D\Theta_\rho=\sum_\tau a_{\rho\tau}\Theta_\tau,
 \qquad a_{\rho\tau}=\langle L_D\Theta_\rho,
 \Theta_\tau\rangle,
\end{equation}
with each row and column finite.  The characters on either side
are orthonormal bases in their fixed-central-character spaces.

\begin{proposition}[\Etag\ Integral-lattice criterion, Part III]\label{prop:int}
The transform $L_D$ permutes the irreducible characters with
positive signs if and only if every $a_{\rho\tau}$ in
\eqref{eq:matrixchars} is an integer.
\end{proposition}
\begin{proof}
An integral unit vector in $\ell^2$ is a signed basis vector;
unitarity gives injectivity and surjectivity because the inverse
is the adjoint of a unitary map.  The value at identity in
\eqref{eq:cD} is positive and fixes the sign.
This is Part III, Corollary 5.2, with its converse.
\end{proof}

\begin{remark}[The actual obstruction]
For wild $v$, Part III does not prove
$a_{\rho\tau}\in\Z$.  The tame Adams argument works with a
finite quotient $Q$ and a central subgroup $C$, under the
hypothesis $\gcd(n,|Q/C|)=1$ (Part III, Lemma 7.5).
At a wild place that hypothesis need not hold.  For the
simplest illustration take $Q=C_p$, $C=1$, $n=p$: then the
Adams operation sends every linear character to the trivial
character.  This does not prove a nonintegral coefficient for
$L_D$; it only shows why that particular proof method fails.
\end{remark}

\subsection{A concrete first wild test}\label{subsec:examplewild}
Consider again $r=2,n=3$, $F=\mathbb Q_3(\zeta_3)$.
Since $[F:\mathbb Q_3]=2$, one has
$|3|_F=3^{-2}$ and
\begin{equation}\label{eq:wildconstants}
 I_F=9,\qquad b_D=3,\qquad c_D=27.
\end{equation}
For the trivial $\rho$ on $D^\times$, compatible with a suitable
choice of genuine central character, $L_D\Theta_\rho$ is
\begin{equation}\label{eq:wildexplicit}
 (L_D1)(z\,t^*)
 =3\,\widetilde\theta(z)
 \frac{|\Nrd(t)|_F}
 {|\operatorname{Trd}(t)^2-\Nrd(t)|_F}
\end{equation}
for noncentral $t$ (and zero on nongood regular lifts).
Indeed for the eigenvalues $\lambda_1,\lambda_2$ one has
\[
 \frac{\Delta(t)}{\Delta(t^3)}
 =\left|\frac{\lambda_1\lambda_2}
 {\lambda_1^2+\lambda_1\lambda_2+\lambda_2^2}\right|_F
 =\frac{|\Nrd(t)|_F}{|\operatorname{Trd}(t)^2-\Nrd(t)|_F}.
\]
Near a scalar $t=a$, the denominator is $|3a^2|_F$, so
this ratio is $9$, giving $27$ at identity.
Thus any genuine irreducible character equal to this function
must have degree $27$; \emph{every} irreducible genuine character
would have degree divisible by $27$ if the entire transform were a
bijection.  Neither the existence of such a $27$-dimensional
representation nor the absence of smaller genuine representations
has been established here.

A finite-level test would select a sufficiently small
conjugation-invariant open subgroup on which the relevant
representations are trivial (compatible with their central
character), form the corresponding \emph{twisted} finite quotient
of $H/Z$, calculate its genuine or projective character table,
and evaluate the finite inner products \eqref{eq:matrixchars}.
A plain quotient $\tH/(\tZ U)$ cannot be used by itself as a
``genuine character table'': dividing out by $\tZ$ destroys the
prescribed genuine central action.  The central character must be
retained as a multiplier in the finite twisted group algebra.
This computation has not been carried out in this draft.

\section{Connection with Zou's corrected character relation}\label{sec:Zou}
For the Kazhdan--Patterson cover of the \emph{split group},
Zou \cite[\S3]{Zou24} revisits the Flicker--Kazhdan metaplectic
correspondence and records a corrected scalar in the character
relation.  This is an essential input, but its objects are
$\widetilde{\GL_r(F)}$ and ordinary $\GL_r(F)$, not directly
$\tH$ and $H$ at a division place.

\begin{proposition}[\Etag\ Zou's root-sum constant in our case]\label{prop:Zouconstant}
Write $d_F(m)=m|m|_F^{-1/2}$ and
$d_r=\gcd(n,2rc_v+r-1)$.  Under \eqref{eq:parameters},
$d_r=n$.  In Zou's equation (3.2), the coefficient preceding
the root sum in rank $r$ is
\begin{equation}\label{eq:Zouprefactor}
 \frac{1}{d_F(d_r)|n|_F^{r/2}}
 =\frac{1}{n|n|_F^{(r-1)/2}}=\frac{b_D}{n}.
\end{equation}
For a regular elliptic element $x\in\GL_r(F)$ with field
centralizer of degree $r$, the equation $t^n=x^n$ has exactly
$n$ solutions $t=\zeta x$, $\zeta\in\mu_n$, provided $x^n$ is
regular.  If the ordinary linear character is trivial on the
scalar $\mu_n$, all those terms agree, and the right-hand side
therefore has the effective coefficient $b_D$.
\end{proposition}
\begin{proof}
The equation $d_r=n$ is immediate from scalar centrality;
\eqref{eq:Zouprefactor} follows by substitution.  Any solution
$t$ centralizes $x^n$, so $t\in F[x]^\times$ because $x^n$
is regular elliptic.  In the field $F[x]$, $(t/x)^n=1$,
and $\mu_n(F[x])=\mu_n(F)$.  Thus $t=\zeta x$.
The Weyl discriminant is invariant under central scalar
multiplication, the ordinary character takes the same value
because it is trivial on $\mu_n$, and $(\zeta x)^*=x^*$
by the explicit odd-$n$ scalar star computation in
Section~\ref{subsec:star}.
These are precisely the equal terms in Zou's (3.2).
\end{proof}

\begin{remark}[What is and is not composed]
At the ramified places one may compare the \emph{desired}
$\GDb_v\leftrightarrow\Gb_v$ character relation of Parts I--II
with (a) the classical linear Jacquet--Langlands elliptic
character identity and (b) Zou's formula.  The common coefficient
$b_D$ and the familiar sign $(-1)^{r-1}$ are consistent with
this comparison.  Nevertheless \emph{the existence of an
irreducible genuine image on the covered division side is the
missing conclusion}; combining two known formulas cannot be
used to assume that missing representation exists.
For non-elliptic split centralizers, retain Zou's entire
root sum, including the central-character factors.
\end{remark}

\section{Global geometric comparison: proved input and new obligations}\label{sec:geometric}
There are two comparisons to keep distinct:
\begin{equation}\label{eq:twocomparisons}
 \begin{gathered}
 \GDb(\A_K)\ \longleftrightarrow\ \Gb(\A_K)
 \quad\text{(Parts I--II, under their hypotheses)},\\[3pt]
 D^\times(\A_K)\ \dashrightarrow\ \GDb(\A_K)
 \quad\text{(proposed additional comparison)}.
 \end{gathered}
\end{equation}
The first has a proved simple trace-formula comparison with
regular-support or later matching hypotheses (Part I, Theorem 6.6;
Part II, Theorem 6.2, in the numbering cited by the referee).
The second is not a formal consequence of the first or of the
unitarity of $L_D$.

\subsection{Geometric coefficients and the product formula}
For a regular elliptic rational $\gamma$, with $E=K[\gamma]$,
Part I identifies the coefficient
\[
 c(\gamma)=\frac{\vol(\A_K^\times E^\times\backslash
 \A_E^\times;dt)}{|W_\gamma|}.
\]
The same torus measure $dt$ must be used whenever this term
is compared with the corresponding linear inner-form orbital integral.
Part I uses the normalized orbital integral
$J_v(\gamma,f)=\Delta_v(\gamma)O_v(\gamma,f)$.
For a rational regular element the product formula implies
\begin{equation}\label{eq:discprod}
 \prod_v\Delta_v(\gamma)=1.
\end{equation}
This is because each discriminant is the normalized absolute
value of a nonzero element of $K$ (or its square root), including
the determinant normalization in \eqref{eq:Delta}.
The local \emph{ratio} of discriminants
$\Delta_v(t)/\Delta_v(t^n)$ must be kept until local transfer
and the global product are both defined.

At $v\in S$ Part III, Theorem 9.1, gives the exact local
normalized orbital-integral identity
\begin{equation}\label{eq:localJO}
 J_{\tH_v}\bigl(z\,t^*,\mathcal O_{D,v}f_v\bigr)
 =b_{D,v}\,\widetilde\theta_v(z)^{-1}
 J_{H_v}(t,f_v),
\end{equation}
using the same quotient measures on the two copies of $H_v/Z_v$.
At $v\notin S$, $H_v\cong\GL_r(K_v)$ and the appropriate
root-sum comparison is not \eqref{eq:localJO}; it must use
Flicker--Kazhdan--Zou, and it can involve the kernel
$\mu_n^{k_v}/\mu_n$ of Proposition \ref{prop:power-etale}.
In particular the local factors $b_{D,v}$ for $v\in S$
\emph{cannot} simply be cancelled using
$\prod_v|n|_v=1$ without the split-side contributions.

\subsection{Haar ledger}\label{subsec:ledger}
For $v\in S$, Part II, Definition 2.1, fixes
\begin{equation}\label{eq:formaldegree}
 \vol(D_v^\times/F_v^\times;dg_{D,v})
 =d(\mathrm{St};dg_{\GL,v})^{-1},
\end{equation}
where $d(\mathrm{St};dg_{\GL,v})$ is the formal degree
relative to the chosen $\GL_r(F_v)/F_v^\times$ measure.
The probability measure used in Part III is
$d(\mathrm{St};dg_{\GL,v})\,dg_{D,v}$.
When \emph{both} sides of \eqref{eq:localJO} are scaled by the
same measure, $b_{D,v}$ is unchanged.  Different measures on the
$\GL$ side must still be reconciled with \eqref{eq:formaldegree}.

\begin{center}\small
\begin{longtable}{@{}p{0.27\linewidth}p{0.53\linewidth}p{0.10\linewidth}@{}}
\toprule
Item & Required statement or calculation & Status\\\midrule
\endhead
Ambient cocycle and common center & Part I, Lemma 5.1 and Prop. 5.5; restriction \eqref{eq:before} & [E]\\
Regular rational good locus & Theorem \ref{thm:globalgood} & [P]\\
Global rational/stable classes & Hilbert 90 and Kottwitz's framework; local embeddability into $D$ & [E]\\
Global torus coefficient & Same $dt$ in \eqref{eq:globalvol}, including $|W_\gamma|$ & [E]\\
Local $v\in S$ orbital transfer & Part III, Theorem 9.1, equation \eqref{eq:localJO} & [E]\\
Local $v\notin S$ root multiplicities & Proposition \ref{prop:power-etale}; matching via [FK], [Zou] & [P]/[O]\\
Unramified unit elements & [FK, \S12] when $(n,N_r)=1$, plus matching of splittings & [E]/[O]\\
Identity and singular terms & Part I, Prop. 6.1; regular-support elimination; general new comparison still needed & [E]/[O]\\
Spectral expansion & Part I--II packet formula with restrictions; new factor matching & [E]/[O]\\
\bottomrule
\end{longtable}
\end{center}
Here $N_r$ is the least common multiple of the \emph{composite}
positive integers at most $r$ (the convention in [FK]).  The
classical unramified matching used in [FK, \S12] is established
under $(n,N_r)=1$; outside that range one must supply an alternative
theorem rather than erase the assumption.

\subsection{A precise conditional geometric assertion}
\begin{proposition}[\Ctag\ Termwise trace comparison criterion]\label{prop:conditionalglobal}
Choose matching factorizable test functions, with a regular elliptic
support condition at one place eliminating the identity and singular
classes, and a cuspidal place whenever the chosen simple trace
formula requires it.  Suppose that:
\begin{enumerate}[label=(\roman*)]
\item the covered and ordinary test functions are normalized at every
place, including unramified places, and all local root sums are included;
\item for every globally good regular rational class $[s^n]$ the product
of the local normalized orbital integrals, after any prescribed
local root summations, agrees \emph{with its central-character
phase} on the two sides;
\item the same torus Haar measure is used, the global rational
splittings are compatible, and the relevant trace expansions are
absolutely convergent.
\end{enumerate}
Then the geometric sides of the two simple trace formulas agree,
and so do their spectral sides.
\end{proposition}
\begin{proof}
By Theorem \ref{thm:globalgood} the noncentral rational
summation sets agree, and by \eqref{eq:Wstabilizer} their normalizer
stabilizer orders agree.  Kottwitz's induced-torus and Brauer-group
conditions determine the common represented classes.  The torus
volume factors are equal by (iii); discriminants and local
orbital-integral products match by (i)--(ii) and
\eqref{eq:discprod}.  The support assumption suppresses the
identity/singular terms.  Absolute convergence permits termwise
comparison, proving equality of geometric expansions, and hence
of their spectral expansions.  This is a formal \emph{conditional}
reduction: (i)--(iii) encode substantive as-yet-unproved local and
global compatibilities for the new comparison.
\end{proof}

\subsection{Why Li, Kottwitz, and DKV do not close the gap}\label{subsec:litroute}
The linear Jacquet--Langlands theorem of
Deligne--Kazhdan--Vign\'eras \cite{DKV} controls the ordinary
$D^\times$--$\GL_r$ relation, including the elliptic character
sign, but it does not automatically manufacture a genuine
representation of $\GDb_v$.  Kottwitz's analysis of stable
conjugacy and relative trace coefficients \cite{Kottwitz84}
provides the correct arithmetic setting; for the induced tori
$\operatorname{Res}_{E/K}\mathbb G_m$ one has
$H^1(K,T)=1$ by Shapiro's lemma and Hilbert 90, and embeddings
$E\hookrightarrow D$ are controlled by Brauer splitting.
Wen--Wei Li proves an invariant trace formula for appropriate
adelic covering groups, under the trace Paley--Wiener hypothesis
at real archimedean places \cite{Li14}.  This hypothesis is known
for the cited metaplectic covers of $\GL$, but Li's general theorem
does not assert the particular \emph{equality} required in
\eqref{eq:twocomparisons}.  The simpler already-proved trace formulas
of Parts I--II remain the appropriate initial comparison.

\section{Spectral separation and a research program}\label{sec:spectral}
\subsection{The actual spectral distribution identity}
Let $H=D^\times$ and $\widetilde H=\GDb$ globally, with
compatible central characters $\theta,\widetilde\theta$.
The quotient $H(K)Z_H(\A_K)\backslash H(\A_K)$ is compact
because $D$ is a division algebra over $K$; its covered analogue
is compact modulo the specified central preimage as well.
Thus the fixed-central-character right regular representations
have discrete Hilbert decompositions.  For factorizable test
functions with the inverse central transformation law, define
\begin{align}
 I_H(f)&=\sum_{\rho\in\mathcal A(H,\theta)}
       m_H(\rho)\prod_v\operatorname{tr}\rho_v(f_v),\\
 I_{\widetilde H}(\widetilde f)
     &=\sum_{\widetilde\rho\in\mathcal A(\widetilde H,\widetilde\theta)}
       m_{\widetilde H}(\widetilde\rho)
       \prod_v\operatorname{tr}\widetilde\rho_v(\widetilde f_v).
       \label{eq:spectralglobal}
\end{align}
Here $\mathcal A$ denotes the relevant automorphic spectrum
and $m$ is its actual spectral multiplicity, not an assumed $1$.
The expressions are interpreted on trace-class test functions
satisfying the simple trace-formula hypotheses; a bare
unrestricted infinite sum is not justified by notation alone.
Under the assumptions of Proposition~\ref{prop:conditionalglobal},
the equality of the geometric expansions supplies
$I_H(f)=I_{\widetilde H}(\widetilde f)$.  It does not
identify summands.  This is a conditional identity of
\emph{distributions}, with both the multiplicities and local
trace factors retained.

\subsection{Individual local characters and integrality}
An equality of global spectral sums is a statement about
\emph{distributions}.  To isolate a single local irreducible at a
wild place, one needs an argument that the relevant coefficients
are integers (or multiplicities which can be shown to isolate a
single constituent).  Linear independence alone can separate
finite lists of characters; it does not supply the required
integrality from a unitary Fourier transform.

\begin{proposition}[\Ctag\ From spectral coefficients to one character]
Fix $v\in S$ wild and an irreducible ordinary character
$\Theta_\rho$ of $D_v^\times$.  If a fully justified spectral
comparison proves that each coefficient $a_{\rho\tau}$ in
\eqref{eq:matrixchars} is an integer, then
$L_{D,v}\Theta_\rho$ is a single irreducible genuine character.
If this is true for every $\rho$, the map is a bijection on the
fixed-central-character irreducibles.
\end{proposition}
\begin{proof}
Apply Proposition \ref{prop:int}.  The entire force of the
hypothesis is in the integer-coefficient assertion; neither the
preceding geometric criterion nor $L^2$ unitarity alone gives it.
\end{proof}

\paragraph{Separating the remaining tasks.}
First, the local wild problem can be tested using the explicit
rank-two example in Section \ref{subsec:examplewild}; it requires
an actual twisted finite character table and a proof that the
chosen finite level captures the desired row.  Second, at split
places the root sum in Zou's (3.2) should be written against the
kernel and index in Proposition \ref{prop:power-etale}, including
the chosen Haar measures and central phases.  Third, one needs
matching unramified test functions under the precise
$(n,N_r)=1$ or an improved hypothesis, and a complete global
coefficient identity as in Proposition \ref{prop:conditionalglobal}.
Finally, one must justify globalization and spectral separation
for representations of the covered inner form, and remove the
auxiliary supercuspidal and non-cancellation hypotheses of Parts
I--II.  These are mathematical work packages, not restatements of
an already-proved theorem.

\section{Selected glossary}\label{sec:glossary}
\begin{description}[style=nextline,leftmargin=3.0cm]
\item[Good element] A regular lift commuting with the full preimage of its ordinary centralizer; genuine invariant characters vanish on the nongood regular locus.
\item[Scalar-central cover] A cover whose entire inverse image of $F^\times I_r$ is central; here forced by $n\mid r-1+2rc_v$.
\item[Power-map kernel] $\ker(P_n)$ on $T(F)/F^\times$, distinct from the image index.  The kernel can be nontrivial on product tori even when the corresponding field-torus map is injective.
\item[Star map] $h^*=\widetilde h^{\,n}$, independent of lift; multiplication on tori uses $n$ odd.
\item[Weyl discriminant] The factor $\Delta=|\det(1-\operatorname{Ad}(t)|\mathfrak g/\mathfrak t)|^{1/2}$; it is not the unnormalized characteristic-polynomial discriminant.
\item[Twisted finite quotient] A finite quotient of $H/Z$ with the central character retained as a multiplier, used to compute projective/genuine irreducible characters.
\item[Packet identity] An identity of traces over collections of automorphic representations; it does not assert a bijection between individual local irreducibles.
\item[Formal degree] The Plancherel normalization of a square-integrable representation relative to a specified quotient Haar measure.
\item[Root-sum relation] Zou's equation (3.2), which sums over all $n$th roots satisfying the stated regularity condition rather than choosing one arbitrarily.
\end{description}

\section*{Acknowledgments}
This research note was prepared with mathematical and editorial
assistance from ChatGPT (OpenAI).  Claude (Anthropic) supplied the
referee report on Draft 1, including identification of the existing
Flicker--Kazhdan nonsplit-torus theorem and suggested repairs to
the chronology of the series.  Earlier Parts I--III were developed
with separate assistance acknowledged there.  As with any
mathematical research draft, statements labeled conditional or open
require independent verification or further work before a final
correspondence theorem can be claimed.

\begin{thebibliography}{99}
\bibitem{DKV} P. Deligne, D. Kazhdan and M.-F. Vign\'eras,
\emph{Repr\'esentations des alg\`ebres centrales simples $p$-adiques},
in \emph{Repr\'esentations des groupes r\'eductifs sur un corps local},
Travaux en Cours 8, Hermann, Paris (1984), 33--117.
\bibitem{FK} Y. Flicker and D. Kazhdan,
\emph{Metaplectic correspondence}, Publ. Math. IH\'ES \textbf{64}
(1986), 53--110, especially \S\S2--4, 12, 18, 24--28.
\url{https://www.numdam.org/articles/10.1007/BF02699192/}.
\bibitem{Zou24} J. Zou,
\emph{Local metaplectic correspondence and applications},
arXiv:2209.05852v3 (19 December 2024), especially
(2.2), (3.2), Remarks 3.1 and 3.8, and Proposition 3.6.
\url{https://arxiv.org/abs/2209.05852}.
\bibitem{KP} D. Kazhdan and S. Patterson,
\emph{Metaplectic forms}, Publ. Math. IH\'ES \textbf{59}
(1984), 35--142.
\bibitem{Kottwitz84} R. E. Kottwitz,
\emph{Stable trace formula: cuspidal tempered terms}, Duke Math. J.
\textbf{51} (1984), 611--650.
\bibitem{Li14} W.-W. Li,
\emph{La formule des traces pour les rev\^etements de groupes
r\'eductifs connexes. IV. Distributions invariantes},
Ann. Inst. Fourier \textbf{64} (2014), 2379--2448.
\url{https://doi.org/10.5802/aif.2915}.
\bibitem{Takeda} S. Takeda,
\emph{Remarks on metaplectic tensor products for covers of $\GL_r$},
Pacific J. Math. \textbf{290} (2017), 415--449.
\bibitem{JoynerI} D. Joyner with collaborators,
\emph{A metaplectic Jacquet--Langlands transfer for division algebras
of prime degree, I}, Drafts 5--6 (2026), especially \S\S2--6.
\bibitem{JoynerII} D. Joyner with collaborators,
\emph{A metaplectic Jacquet--Langlands transfer for division algebras
of prime degree, II}, Drafts 4--6 (2026), especially \S\S2, 5--7.
\bibitem{JoynerIII} D. Joyner with collaborators,
\emph{Toward a metaplectic Jacquet--Langlands transfer, III:
The division-algebra side and its local character transform},
Drafts 6--7 (2026), especially \S\S3--5, 7--9, 11--12.
\end{thebibliography}
\end{document}
